\title{DeepSeekMath: Advancing the Boundaries of Mathematical Reasoning in Open Language Models}

\begin{abstract}
Mathematical reasoning represents a formidable challenge for language models due to its intricate and organized structure. In this study, we present DeepSeekMath~7B, an extension of DeepSeek-Coder-Base-v1.5 7B, further pre-trained on 120B math-centric tokens extracted from Common Crawl, alongside natural language and coding datasets. DeepSeekMath~7B attains a remarkable 51.7\% accuracy on the competition-level MATH benchmark without the use of external tools or ensemble methods, nearing the capabilities of Gemini-Ultra and GPT-4. Employing self-consistency across 64 samples, DeepSeekMath~7B reaches 60.9\% on MATH. The model’s proficiency in mathematical reasoning is driven by two principal elements: firstly, leveraging the vast potential of publicly accessible web data through a carefully designed data selection framework; secondly, the introduction of Group Relative Policy Optimization (GRPO), a PPO variant that simultaneously boosts mathematical reasoning performance and optimizes PPO’s memory consumption.
\end{abstract}

\section{Introduction}

Large language models (LLMs) have transformed the methodology for mathematical reasoning within artificial intelligence, catalyzing notable progress in both the quantitative reasoning benchmark \citep{MATH} and the geometric reasoning benchmark \citep{trinh2024solving}. Additionally, these models have demonstrated effectiveness in aiding humans to tackle complex mathematical challenges \citep{tao}. Nonetheless, state-of-the-art models like GPT-4 \citep{gpt4} and Gemini-Ultra \citep{gemini} remain inaccessible to the public, and currently available open-source alternatives lag significantly behind in terms of performance.

In this work, we present DeepSeekMath, a domain-specific language model that markedly surpasses the mathematical performance of open-source models and nears the benchmark results of GPT-4 on academic tests. To accomplish this, we develop the DeepSeekMath Corpus, a large-scale, high-quality pre-training dataset containing 120B math-related tokens. This dataset is derived from the Common Crawl (CC) by utilizing a fastText-based classifier \citep{joulin2016fasttext}. Initially, the classifier is trained using positive samples from OpenWebMath \citep{openwebmath} alongside a diverse array of other web pages as negative samples. Following this, the classifier is applied to mine further positive examples from the CC, which are then refined through human annotation. The classifier is subsequently retrained with this enriched dataset to enhance its accuracy. Evaluation results demonstrate that the extensive corpus is of excellent quality, as evidenced by our base model DeepSeekMath-Base 7B achieving 64.2\% on GSM8K \citep{gsm8k} and 36.2\% on the challenging MATH dataset \citep{MATH}, outperforming Minerva 540B \citep{minerva}. Furthermore, since the DeepSeekMath Corpus is multilingual, improvements are also observed on Chinese mathematical benchmarks \citep{wei2023cmath,agieval}. We consider our methodology in mathematical data processing as an initial framework for the research community, with substantial potential for future advancements.

DeepSeekMath-Base is initialized from DeepSeek-Coder-Base-v1.5 7B \citep{deepseek-coder}, as we find that beginning with a code-trained model is a superior option compared to using a general LLM. Additionally, we observe that mathematical training also enhances the model’s performance on the MMLU \citep{mmlu} and BBH benchmarks \citep{bbh}, suggesting that it not only strengthens the model’s mathematical skills but also improves its overall reasoning abilities.

Following pre-training, we perform mathematical instruction tuning on DeepSeekMath-Base using data incorporating chain-of-thought \citep{cot}, program-of-thought \citep{pot,pal}, and tool-integrated reasoning \citep{tora}. The resulting model, DeepSeekMath-Instruct 7B, surpasses all other 7B models and is competitive with open-source instruction-tuned models at the 70B scale.

Moreover, we propose Group Relative Policy Optimization (GRPO), a variant of the Proximal Policy Optimization (PPO) reinforcement learning (RL) algorithm \citep{schulman2017proximal}. GRPO eliminates the need for a critic model by estimating the baseline from group scores, which greatly reduces the resources required for training. Utilizing only a portion of the English instruction tuning dataset, GRPO achieves significant improvements over the robust DeepSeekMath-Instruct model, both on in-domain tasks (GSM8K: 82.9\% $\rightarrow$ 88.2\%, MATH: 46.8\% $\rightarrow$ 51.7\%) and out-of-domain mathematical challenges (e.g., CMATH: 84.6\% $\rightarrow$ 88.8\%) during the reinforcement learning phase. We also present a unified framework to interpret various approaches such as Rejection Sampling Fine-Tuning (RFT) \citep{yuan2023scaling}, Direct Preference Optimization (DPO) \citep{dpo}, PPO, and GRPO. Within this framework, we find that these methods can all be viewed as either direct or simplified forms of RL. We conduct extensive experiments, including comparisons between online and offline training, outcome versus process supervision, and single-turn versus iterative RL, to thoroughly explore the fundamental components of this paradigm. Finally, we clarify why our RL approach enhances the performance of instruction-tuned models and outline potential future directions for achieving more effective RL based on this unified framework.

\subsection{Contributions}
Our work presents scalable mathematics pre-training along with an in-depth investigation and evaluation of reinforcement learning.

\textbf{Large-Scale Math Pre-Training}

\begin{itemize}[topsep=0pt]
    \item We demonstrate that publicly available Common Crawl data holds significant value for mathematical applications. By employing a carefully crafted data selection process, we create the DeepSeekMath Corpus, a high-quality dataset comprising 120B tokens extracted from web pages filtered specifically for mathematical content. This dataset is nearly 7 times larger than the math web pages used by Minerva \citep{minerva} and 9 times larger than the recently introduced OpenWebMath \citep{openwebmath}.

    \item Our base model, DeepSeekMath-Base 7B, pre-trained on this data, achieves performance on par with Minerva 540B \citep{minerva}, suggesting that the sheer number of parameters is not the sole determinant of mathematical reasoning capabilities. Smaller models trained on high-quality datasets can also deliver strong results.

    \item We report insights from our math training experiments. Pre-training on code before math training enhances the model’s proficiency in solving mathematical problems, both with and without the assistance of tools. This provides partial evidence addressing the longstanding question: \textit{does code training enhance reasoning skills?} We contend that it does, at least in the context of mathematical reasoning.

    \item While training on arXiv papers is a common practice, especially in many mathematics-related studies, it yields no significant improvements across the mathematical benchmarks evaluated in this work.
\end{itemize}

\textbf{Investigation and Evaluation of Reinforcement Learning}

\begin{itemize}[topsep=0pt]
    \item We propose Group Relative Policy Optimization (GRPO), a reinforcement learning algorithm that is both efficient and effective. GRPO eliminates the need for a critic model by estimating the baseline from group scores, which drastically lowers the training resource requirements compared to Proximal Policy Optimization (PPO).
    \item We show that GRPO substantially improves the performance of our instruction-tuned model, DeepSeekMath-Instruct, by utilizing only instruction-tuning data. Additionally, we observe improvements in out-of-domain performance during the reinforcement learning stage.
    \item We offer a unified framework to interpret various approaches, including RFT, DPO, PPO, and GRPO. We also perform extensive experiments—such as comparing online versus offline training, outcome versus process supervision, and single-turn versus iterative reinforcement learning—to thoroughly explore the core components of this framework.
    \item Using our unified framework, we investigate the factors behind the success of reinforcement learning and outline several promising avenues for achieving more effective reinforcement learning of large language models (LLMs).
\end{itemize}

\subsection{Summary of Evaluations and Metrics}

\begin{itemize}[topsep=0pt]
    \item \textbf{English and Chinese Mathematical Reasoning}: We conduct thorough evaluations of our models on English and Chinese benchmarks, encompassing math problems ranging from elementary school to college level. English datasets include GSM8K \citep{gsm8k}, MATH \citep{MATH}, SAT \citep{llemma}, OCW Courses \citep{minerva}, and MMLU-STEM \citep{mmlu}. Chinese datasets include MGSM-zh \citep{mgsm}, CMATH \citep{wei2023cmath}, Gaokao-MathCloze \citep{agieval}, and Gaokao-MathQA \citep{agieval}. We assess the models’ capacity to generate self-contained textual solutions without external tools, as well as their ability to solve problems by utilizing Python.
    
    On English benchmarks, DeepSeekMath-Base achieves competitive results compared to the closed-source Minerva 540B \citep{minerva} and outperforms all open-source base models (e.g., Mistral 7B \citep{mistral} and Llemma-34B \citep{llemma}), whether or not they have undergone math-specific pre-training, often by a notable margin. Remarkably, DeepSeekMath-Base excels on Chinese benchmarks, likely because our approach diverges from prior work \citep{minerva,llemma} that gathered math pre-training data exclusively in English, and instead incorporates high-quality non-English data. With the addition of mathematical instruction tuning and reinforcement learning, DeepSeekMath-Instruct and DeepSeekMath-RL achieve strong results, surpassing 50\% accuracy on the competition-level MATH dataset for the first time in the open-source domain.
\end{itemize}

\begin{itemize}
    \item \textbf{Formal Mathematics}: We assess DeepSeekMath-Base using the informal-to-formal theorem proving task described in \citep{dsp_proof} on miniF2F \citep{minif2f}, selecting Isabelle \citep{isabelle} as the proof assistant. DeepSeekMath-Base exhibits robust few-shot autoformalization capabilities.
    \item \textbf{Natural Language Understanding, Reasoning, and Code}: To develop a thorough understanding of the models' general comprehension, reasoning, and coding abilities, we test DeepSeekMath-Base on the Massive Multitask Language Understanding (MMLU) benchmark \citep{mmlu}, which includes 57 multiple-choice tasks spanning various domains; BIG-Bench Hard (BBH) \citep{bbh}, comprising 23 demanding tasks that largely require multi-step reasoning; and HumanEval \citep{codex} and MBPP \citep{mbpp}, which are commonly used to evaluate code language models. Pre-training on mathematics enhances both language understanding and reasoning performance.
\end{itemize}

\section{Math Pre-Training}

\subsection{Data Collection and Decontamination} In this section, we detail the procedure for creating the DeepSeekMath Corpus from Common Crawl. As illustrated in Figure \ref{fig:data_collect}, we introduce an iterative pipeline that systematically collects a large-scale mathematical dataset from Common Crawl, beginning with a seed corpus (for example, a small but high-quality math-related dataset). It should be noted that this methodology can be applied to other areas such as coding.

Initially, we select OpenWebMath \citep{openwebmath}, a high-quality collection of mathematical web texts, as our seed corpus. Using this dataset, we train a fastText model \citep{joulin2016fasttext} to identify additional web pages similar to those in OpenWebMath. Concretely, we randomly pick 500,000 samples from the seed corpus as positive training examples and another 500,000 web pages from Common Crawl as negative examples. We utilize an open-source library\footnote{\url{https://fasttext.cc}} to train the model, setting the vector size to 256, learning rate to 0.1, maximum word n-gram length to 3, minimum word frequency to 3, and number of training epochs to 3. To reduce the original Common Crawl size, we apply URL-based deduplication and near-deduplication, resulting in 40B HTML pages. We then use the fastText model to retrieve mathematical web pages from this deduplicated Common Crawl. To exclude low-quality math content, we rank pages based on their fastText-predicted scores and retain only the highest-ranked ones. The amount of retained data is evaluated through pre-training experiments on the top 40B, 80B, 120B, and 160B tokens. For the initial iteration, we decide to keep the top 40B tokens.

Following the initial round of data collection, many mathematical web pages remain uncollected, primarily because the fastText model is trained on a set of positive samples that lack adequate diversity. To address this, we identify additional mathematical web sources to expand the seed corpus, enabling us to fine-tune the fastText model. Specifically, we first partition the entire Common Crawl dataset into separate domains; a domain is defined as a collection of web pages sharing the same base URL. For each domain, we compute the proportion of web pages collected during the first iteration. Domains in which more than 10\% of the pages were collected are deemed math-related (e.g., \url{mathoverflow.net}).

Next, we manually label the URLs within these math-related domains that correspond to mathematical content (e.g., \url{mathoverflow.net/questions}). Web pages linked to these URLs that were not yet collected will be incorporated into the seed corpus. This strategy allows us to obtain more positive examples and thus train a more effective fastText model capable of retrieving a larger quantity of mathematical data in subsequent iterations. After four rounds of data collection, we accumulate 35.5M mathematical web pages, comprising a total of 120B tokens. In the fourth iteration, we observe that nearly 98\% of the data had already been collected during the third iteration, leading us to terminate further data collection.

To prevent benchmark contamination, we adopt the approach of \cite{deepseek-coder} by filtering out web pages containing questions or answers from English mathematical benchmarks such as GSM8K \citep{gsm8k} and MATH \citep{MATH}, as well as Chinese benchmarks like CMATH \citep{wei2023cmath} and AGIEval \citep{agieval}. The filtering rule is as follows: any text segment containing a 10-gram string that exactly matches any substring from the evaluation benchmarks is excluded from our math training corpus. For benchmark texts shorter than 10 grams but at least 3 grams long, we use exact matching to remove contaminated web pages.

\subsection{Assessing the Quality of the DeepSeekMath~Corpus}

We conduct pre-training experiments to compare the DeepSeekMath~Corpus against recently published math-training corpora:

\begin{itemize}[topsep=0pt]
    \item \textbf{MathPile} \citep{mathpile}: a multi-source corpus (8.9B tokens) compiled from textbooks, Wikipedia, ProofWiki, CommonCrawl, StackExchange, and arXiv, with the majority (over 85\%) originating from arXiv;
    \item \textbf{OpenWebMath} \citep{openwebmath}: CommonCrawl data filtered for mathematical content, totaling 13.6B tokens;
    \item \textbf{Proof-Pile-2} \citep{llemma}: a mathematical corpus comprising OpenWebMath, AlgebraicStack (10.3B tokens of mathematical code), and arXiv papers (28.0B tokens). For experiments involving Proof-Pile-2, we follow \cite{llemma} by adopting an arXiv:Web:Code ratio of 2:4:1.
\end{itemize}

\subsubsection{Training Configuration}
\label{sec:quality-policy}
We fine-tune a general pre-trained language model with 1.3B parameters, based on the same architecture as the DeepSeek LLMs \citep{deepseek-llm}, referred to as DeepSeek-LLM 1.3B. Each mathematical dataset is used to separately train the model for 150B tokens. All training processes utilize the efficient and lightweight HAI-LLM \citep{haillm} training framework. Following the DeepSeek LLMs' training protocol, we employ the AdamW optimizer \citep{adamW} with parameters $\beta_1=0.9$, $\beta_2=0.95$, and a $\mathrm{weight\_decay}=0.1$. The learning rate schedule is multi-step: it ramps up to its peak after 2,000 warm-up steps, then declines to 31.6\% of the peak by 80\% of the training duration, and further reduces to 10.0\% of the peak after 90\% of training completion. The peak learning rate is set to 5.3e-4, and training uses a batch size of 4 million tokens with a context length of 4K.

\subsubsection{Evaluation Outcomes}

\textbf{The DeepSeekMath~Corpus is of superior quality, includes multilingual mathematical content, and is the largest dataset available.}

\begin{itemize}[topsep=0pt]
    \item \textbf{Superior Quality}:
    We assess downstream task performance across 8 mathematical benchmarks using few-shot chain-of-thought prompting \cite{cot}. As shown in Table \ref{tab:corpora_comparison}, models trained on the DeepSeekMath~Corpus consistently outperform others. Figure \ref{fig:corpora_comparisons} illustrates that at 50B tokens (equivalent to one full epoch on Proof-Pile-2), the model trained on DeepSeekMath surpasses Proof-Pile-2, suggesting that the average quality of the DeepSeekMath Corpus is higher.
    \item \textbf{Multilingual Coverage}:
    The DeepSeekMath~Corpus includes datasets in various languages, with English and Chinese being the most prevalent. As evidenced in Table \ref{tab:corpora_comparison}, training on the DeepSeekMath~Corpus improves mathematical reasoning abilities in both English and Chinese. By contrast, existing mathematical corpora, which are mostly English-centric, show limited gains and may even impede performance on Chinese mathematical reasoning tasks.
    \item \textbf{Large Scale}:
    The DeepSeekMath~Corpus is multiple times larger than other existing mathematical datasets. Figure \ref{fig:corpora_comparisons} shows that DeepSeek-LLM 1.3B trained on the DeepSeekMath~Corpus exhibits a steeper learning trajectory and more sustained performance improvements. In comparison, baseline corpora are smaller, have been reused across multiple training passes, and the resulting model's performance plateaus quickly.
\end{itemize}

\subsection{Training and Evaluating DeepSeekMath-Base 7B}

In this subsection, we present DeepSeekMath-Base 7B, a foundational model exhibiting strong reasoning capabilities, particularly in mathematics. The model is initialized from DeepSeek-Coder-Base-v1.5 7B \citep{deepseek-coder} and trained on 500 billion tokens. The data composition is as follows: 56\% originates from the DeepSeekMath Corpus, 4\% from AlgebraicStack, 10\% from arXiv, 20\% from Github code, and the remaining 10\% consists of natural language data from Common Crawl in both English and Chinese. We largely follow the training configuration described in Section \ref{sec:quality-policy}, with the exception of setting the maximum learning rate to 4.2e-4 and utilizing a batch size of 10 million tokens.

We perform an extensive evaluation of DeepSeekMath-Base 7B’s mathematical abilities, emphasizing its capacity to generate self-contained mathematical solutions without external aids, solve math problems using tools, and carry out formal theorem proving. In addition to its mathematical proficiency, we provide an overview of the base model’s broader capabilities, including natural language understanding, reasoning, and programming skills.

\paragraph{Mathematical Problem Solving with Step-by-Step Reasoning}

We assess DeepSeekMath-Base's ability to solve mathematical problems using few-shot chain-of-thought prompting \citep{cot} across eight benchmarks in both English and Chinese. These benchmarks cover quantitative reasoning tasks (such as GSM8K \citep{gsm8k}, MATH \citep{MATH}, and CMATH \citep{wei2023cmath}) and multiple-choice problems (including MMLU-STEM \citep{mmlu} and Gaokao-MathQA \citep{agieval}), spanning a wide range of mathematical domains from elementary to college-level difficulty.

As indicated in Table \ref{tab:base_math_cot}, DeepSeekMath-Base 7B achieves the highest performance across all eight benchmarks among open-source base models (including the widely-used general model Mistral 7B \citep{mistral} and the recently introduced Llemma 34B \citep{llemma}, which has been trained on math data from Proof-Pile-2 \citep{llemma}). Remarkably, on the challenging MATH dataset, DeepSeekMath-Base surpasses existing open-source base models by more than 10\% absolute and even exceeds the performance of Minerva 540B \citep{minerva}, a proprietary base model 77 times larger that builds upon PaLM \citep{palm} and is further trained on mathematical literature.

\paragraph{Mathematical Problem Solving with Tool Use} We assess program-assisted mathematical reasoning on GSM8K and MATH via few-shot program-of-thought prompting \citep{pot,pal}. Models are instructed to solve each problem by generating a Python program that can employ libraries like \textit{math} and \textit{sympy} for complex computations. The program’s execution output is then considered the solution. As detailed in Table \ref{tab:base_math_pot_proof}, DeepSeekMath-Base 7B outperforms the previous state-of-the-art model Llemma 34B.

\paragraph{Formal Mathematics}

Automating formal proofs is valuable for ensuring the correctness and dependability of mathematical arguments while improving efficiency, a topic that has garnered increased interest recently. We evaluate DeepSeekMath-Base 7B on the informal-to-formal proving task from \citep{dsp_proof}, which involves generating a formal proof based on an informal statement, its formal equivalent, and an informal proof. The evaluation is conducted on miniF2F \citep{minif2f}, a formal benchmark for Olympiad-level mathematics, where the model generates formal proofs in Isabelle using few-shot prompting. Following \cite{dsp_proof}, the model produces proof sketches that are completed by invoking the automated theorem prover Sledgehammer \citep{sledgehammer}. As shown in Table \ref{tab:base_math_pot_proof}, DeepSeekMath-Base 7B achieves robust results in proof autoformalization.

\paragraph{Natural Language Understanding, Reasoning, and Code}

We measure model capabilities in natural language understanding on MMLU \citep{mmlu}, reasoning on BBH \citep{bbh}, and coding on HumanEval \citep{codex} and MBPP \citep{mbpp}. As reported in Table \ref{tab:understanding_reasoning_code}, DeepSeekMath-Base 7B shows marked improvements in MMLU and BBH performance compared to its predecessor, DeepSeek-Coder-Base-v1.5 \citep{deepseek-coder}, highlighting the beneficial effects of math-focused training on language comprehension and reasoning. Furthermore, by incorporating code tokens during continual training, DeepSeekMath-Base 7B sustains the coding benchmark performance of DeepSeek-Coder-Base-v1.5. Overall, DeepSeekMath-Base 7B substantially outperforms the general-purpose Mistral 7B \citep{mistral} across the reasoning and coding benchmarks.

\section{Supervised Fine-Tuning}

\subsection{SFT Data Curation}

We compile a mathematical instruction-tuning dataset that includes English and Chinese problems drawn from various mathematical disciplines and exhibiting different levels of difficulty: problems are paired with solutions presented in chain-of-thought (CoT) \citep{cot}, program-of-thought (PoT) \citep{pot,pal}, and tool-enhanced reasoning formats \citep{tora}. The dataset contains a total of 776K training instances.

\begin{itemize}[topsep=0pt]
    \item \textbf{English mathematical datasets}: We provide annotations of GSM8K and MATH problems with tool-assisted solutions, and utilize a portion of MathInstruct \citep{MathInstruct} alongside the Lila-OOD training set \citep{lila}, where problems are solved using CoT or PoT methodologies. Our English dataset encompasses a broad range of mathematical areas such as algebra, probability, number theory, calculus, and geometry.
    \item \textbf{Chinese mathematical datasets}: We gather Chinese K-12 math problems across 76 subtopics, including linear equations, with solutions annotated in both CoT and tool-enhanced reasoning formats.
\end{itemize}

\subsection{Training and Evaluating DeepSeekMath-Instruct 7B}

Here, we describe DeepSeekMath-Instruct 7B, which is fine-tuned on mathematical instructions starting from the DeepSeekMath-Base model. Training examples are concatenated randomly up to a maximum context window of 4K tokens. The model is trained for 500 steps using a batch size of 256 and a fixed learning rate of 5e-5.

We assess the model's mathematical capabilities both without and with the use of external tools, across four quantitative reasoning benchmarks in English and Chinese. Our evaluation compares our model to the leading contemporaneous models:

\begin{itemize}[topsep=0pt]
    \item \textbf{Closed-source models} consist of: (1) the GPT series, with GPT-4 \citep{gpt4} and GPT-4 Code Interpreter~\footnote{\url{https://openai.com/blog/chatgpt-plugins\#\#code-interpreter}} being the most advanced, (2) Gemini Ultra and Pro \citep{gemini}, (3) Inflection-2 \citep{inflection-2}, (4) Grok-1~\footnote{\url{https://x.ai/model-card}}, as well as recently released models from Chinese firms including (5) Baichuan-3~\footnote{\url{https://www.baichuan-ai.com}}, and (6) the latest GLM-4~\footnote{\url{https://open.bigmodel.cn/dev/api\#glm-4}} from the GLM family \citep{glm}.
\end{itemize}

These models serve general purposes and most have been subjected to multiple alignment processes.
\item \textbf{Open-source models} encompass: general-purpose models such as (1) DeepSeek-LLM-Chat 67B \citep{deepseek-llm}, (2) Qwen 72B \citep{qwen}, (3) SeaLLM-v2 7B \citep{seallm}, and (4) ChatGLM3 6B \citep{chatglm3}, along with models enhanced for mathematics including (5) InternLM2-Math 20B~\footnote{\url{https://github.com/InternLM/InternLM-Math}}, which is based on InternLM2 and underwent math-specific training followed by instruction tuning, (6) Math-Shepherd-Mistral 7B, which employs PPO training \citep{schulman2017proximal} on Mistral 7B \citep{mistral} using a process-supervised reward model, (7) the WizardMath series \citep{wizardmath}, which enhances mathematical reasoning in Mistral 7B and Llama-2 70B \citep{llama2} through evolve-instruct (a form of instruction tuning utilizing AI-generated instructions) and PPO training with training data mainly drawn from GSM8K and MATH, (8) MetaMath 70B \citep{metamath}, which fine-tunes Llama-2 70B on an expanded version of GSM8K and MATH, (9) ToRA 34B \cite{tora}, a CodeLlama 34B fine-tuned for tool-integrated mathematical reasoning, and (10) MAmmoTH 70B \citep{MathInstruct}, which is Llama-2 70B instruction-tuned on MathInstruct.
\end{itemize}

As detailed in Table \ref{tab:sft_rl_math}, in the evaluation scenario that prohibits tool usage, DeepSeekMath-Instruct 7B exhibits robust step-by-step reasoning capabilities. Remarkably, on the competitive MATH dataset, our model outperforms all open-source models as well as most proprietary models (e.g., Inflection-2 and Gemini Pro) by a margin of at least 9\% absolute. This superiority holds even when compared to notably larger models (e.g., Qwen 72B) or those that have undergone math-specialized reinforcement learning enhancements (e.g., WizardMath-v1.1 7B). While DeepSeekMath-Instruct competes closely with Chinese proprietary models GLM-4 and Baichuan-3 on MATH, it still falls short of GPT-4 and Gemini Ultra.

In the evaluation setting permitting the combination of natural language reasoning and program-based tool use for problem solving, DeepSeekMath-Instruct 7B approaches an accuracy near 60\% on MATH, outperforming all current open-source models. On other benchmarks, our model performs competitively against DeepSeek-LLM-Chat 67B, the prior state-of-the-art model that is ten times larger.

\section{Reinforcement Learning}

\subsection{Group Relative Policy Optimization} Reinforcement learning (RL) has demonstrated effectiveness in further enhancing the mathematical reasoning capabilities of large language models (LLMs) following the Supervised Fine-Tuning (SFT) phase \citep{wang2023math,wizardmath}. In this section, we present our efficient and effective RL algorithm, Group Relative Policy Optimization (GRPO).

\subsubsection{From PPO to GRPO} Proximal Policy Optimization (PPO) \citep{schulman2017proximal} is a widely used actor-critic RL algorithm in the RL fine-tuning of LLMs \citep{ouyang2022training}. Specifically, it improves LLMs by maximizing the following surrogate objective:

\begin{equation}
\footnotesize
    \mathcal{J}_{PPO}(\theta) = \mathbb{E}_{q \sim P(Q), o \sim \pi_{\theta_{old}}(O|q)} \frac{1}{|o|} \sum_{t=1}^{|o|} \min \left[ \frac{\pi_\theta(o_{t} | q, o_{<t})}{\pi_{\theta_{old}}(o_{t} | q, o_{<t})} A_{t}, \text{clip} \left( \frac{\pi_\theta(o_{t} | q, o_{<t})}{\pi_{\theta_{old}}(o_{t} | q, o_{<t})}, 1 - \varepsilon, 1 + \varepsilon \right) A_{

Here, $\pi_{\theta}$ and $\pi_{\theta_{old}}$ denote the current and previous policy models, respectively, while $q$ and $o$ represent questions and outputs sampled from the question dataset and the old policy $\pi_{\theta_{old}}$. The parameter $\varepsilon$ is a clipping-related hyper-parameter introduced in PPO to stabilize the training process. The advantage $A_t$ is computed using Generalized Advantage Estimation (GAE) \citep{gae}, relying on the rewards $\{r_{\ge t}\}$ and a learned value function $V_{\psi}$. Consequently, in PPO, the value function is trained alongside the policy model. To prevent excessive optimization of the reward model, the common practice is to incorporate a per-token KL penalty from a reference model into the reward for each token \citep{ouyang2022training}, expressed as:

\begin{equation}
    r_{t} = r_\varphi(q, o_{\le t}) - \beta \log\frac{\pi_{\theta}(o_{t}|q, o_{<t})}{\pi_{ref}(o_{t}|q, o_{<t})},
\label{eq:PPO-reward}
\end{equation}

where $r_\varphi$ is the reward model, $\pi_{ref}$ stands for the reference model (typically the initial SFT model), and $\beta$ is the coefficient controlling the KL penalty.

Since the value function used in PPO is often a separate model of comparable size to the policy model, this adds significant memory and computational overhead. Moreover, during RL training, the value function serves as a baseline in advantage calculation to reduce variance. In the context of large language models, the reward model typically assigns a reward only to the final token, which complicates training a value function that accurately evaluates each token. To tackle this, as illustrated in Figure \ref{fig:grpo}, we propose Group Relative Policy Optimization (GRPO), which eliminates the need for an additional value function as required in PPO. Instead, it uses the average reward of multiple sampled outputs generated in response to the same question as the baseline. Specifically, for each question $q$, GRPO samples a set of outputs $\{o_1, o_2, \cdots, o_G\}$ from the old policy $\pi_{\theta_{old}}$ and optimizes the policy model by maximizing the following objective:

\begin{equation}
\footnotesize
\begin{split}
    \mathcal{J}_{GRPO}(\theta) &= \mathbb{E}{[q \sim P(Q), \{o_i\}_{i=1}^G \sim \pi_{\theta_{old}}(O|q)]}  \\
    & \frac{1}{G}\sum_{i=1}^G\frac{1}{|o_i|} \sum_{t=1}^{|o_i|} \left\{ \min \left[ \frac{\pi_\theta(o_{i,t} | q, o_{i,<t})}{\pi_{\theta_{old}}(o_{i,t} | q, o_{i,<t})} \hat{A}_{i,t}, \text{clip} \left( \frac{\pi_\theta(o_{i,t} | q, o_{i,<t})}{\pi_{\theta_{old}}(o_{i,t} | q, o_{i,<t})}, 1 - \varepsilon, 1 + \varepsilon \right)  \hat{A}_{i,t} \right] - \beta \mathbb{D}_{KL}\left[\pi_{\theta} || \pi_{ref}\right]\right\} ,
\end{split}
\label{eq:GRPO-obj}
\end{equation}

where $\varepsilon$ and $\beta$ are hyper-parameters, and $\hat{A}_{i,t}$ denotes the advantage computed solely from the relative rewards of the outputs within each group, to be explained in the subsequent subsections. This group-relative calculation of advantages fits well with the comparative nature of reward models, which are generally trained on datasets comparing outputs for the same question. Note also that instead of incorporating the KL penalty within the reward, GRPO applies regularization by directly adding the KL divergence between the trained policy and the reference policy to the loss function, thereby simplifying the computation of $\hat{A}_{i,t}$. Unlike the KL penalty term in (\ref{eq:PPO-reward}), we estimate the KL divergence using the following unbiased estimator \citep{kl_approx}:

\begin{equation}
\small
    \mathbb{D}_{KL}\left[\pi_{\theta} || \pi_{ref}\right

\begin{algorithm}[t]
  \small
  \caption{Iterative Group Relative Policy Optimization}
  \textbf{Input} initial policy model $\pi_{\theta_{\text{init}}}$; reward models $r_\varphi$; task prompts $\mathcal{D}$; 
  hyperparameters $\varepsilon$, $\beta$, $\mu$
  \begin{algorithmic}[1]
    \State Initialize policy model $\pi_\theta \leftarrow \pi_{\theta_{\text{init}}}$
    \For{iteration = 1, \dots, I}
       \State Set reference model $\pi_{ref} \leftarrow \pi_{\theta}$
      \For{step = 1, \dots, M}
      \State Draw a batch $\mathcal{D}_b$ from $\mathcal{D}$
      \State Assign old policy model $\pi_{\theta_{old}} \leftarrow \pi_{\theta}$ 
      \State Generate $G$ outputs $\{o_i\}_{i=1}^G \sim \pi_{\theta_{old}} (\cdot \mid q) $ for each question $q \in \mathcal{D}_b$
      \State Evaluate rewards $\{r_i\}_{i=1}^{G}$ for each output $o_i$ using $r_{\varphi}$ 
      \State Compute $\hat{A}_{i,t}$ for token $t$ in $o_i$ via group relative advantage estimation.
      \For{GRPO iteration = 1, \dots, $\mu$}
        \State Update policy model $\pi_{\theta}$ by maximizing the GRPO objective (Equation \ref{eq:GC-GRPO})
      \EndFor
    \EndFor \State Continuously update $r_\varphi$ using a replay buffer.
    \EndFor 
  \end{algorithmic}
  \textbf{Output} $\pi_\theta$
  \label{alg:iter-grpo}
\end{algorithm}

\subsubsection{Outcome Supervision RL with GRPO} Formally, for each question $q$, a set of outputs $\{o_1, o_2, \cdots, o_G\}$ is sampled from the previous policy model $\pi_{\theta_{old}}$. A reward model then assigns scores to these outputs, producing $G$ rewards $\mathbf{r}=\{r_1, r_2, \cdots, r_G\}$. These rewards are normalized by subtracting their group mean and dividing by the group standard deviation. Outcome supervision supplies the normalized reward at the conclusion of each output $o_i$ and assigns the advantage values $\hat{A}_{i, t}$ for all tokens in the output as this normalized reward, i.e., $\hat{A}_{i, t} = \widetilde{r}_i = \frac{r_i- {\rm mean}(\mathbf{r})}{{\rm std}(\mathbf{r})}$, then optimizes the policy by maximizing the objective given in equation (\ref{eq:GRPO-obj}).

\subsubsection{Process Supervision RL with GRPO} Since outcome supervision only provides a reward at the end of each output, which might be insufficient or inefficient for guiding the policy in complex mathematical problems, we also investigate process supervision following \cite{wang2023math}, which offers a reward at the conclusion of each reasoning step. Formally, for question $q$ and $G$ sampled outputs $\{o_1, o_2, \cdots, o_G\}$, a process reward model scores each step within the outputs, producing rewards: $\mathbf{R} = \{ \{r_1^{index(1)},\cdots,r_1^{index(K_1)}\}, \cdots,  \{r_G^{index(1)},\cdots,r_G^{index(K_G)}\} \}$, where $index(j)$ denotes the token index at the end of the $j$-th step, and $K_i$ is the total number of steps in the $i$-th output. These rewards are normalized by their overall mean and standard deviation, i.e., $\widetilde{r}_i^{index(j)} = \frac{r_i^{index(j)} - {\rm mean}(\mathbf{R})}{{\rm std}(\mathbf{R})}$. Then, process supervision computes the advantage of each token as the sum of normalized rewards from all subsequent steps, i.e., $\hat{A}_{i, t} = \sum_{index(j) \ge t} \widetilde{r}_i^{index(j)}$, and proceeds to optimize the policy by maximizing the objective in equation

\subsubsection{Iterative RL with GRPO} As the reinforcement learning process advances, the previously trained reward model might no longer adequately guide the current policy model. Hence, we investigate iterative RL using GRPO. As detailed in Algorithm \ref{alg:iter-grpo}, iterative GRPO involves generating new training datasets for the reward model by sampling from the policy model. The existing reward model is continuously updated via a replay strategy that incorporates 10\% of past data. Subsequently, the policy model is set as the reference model and is iteratively trained using the updated reward model.

\subsection{Training and Evaluating DeepSeekMath-RL}

Our RL experiments are based on DeepSeekMath-Instruct 7B. The RL training data consist of chain-of-thought formatted questions from GSM8K and MATH found within the SFT dataset, totaling approximately 144K questions. We exclude other SFT questions to specifically examine RL’s impact on benchmarks with limited data during the RL stage. The reward model training set is constructed in accordance with \citep{wang2023math}. The initial reward model is trained starting from DeepSeekMath-Base 7B with a learning rate of 2e-5. For GRPO, the policy model’s learning rate is set to 1e-6, with a KL divergence coefficient of 0.04. For each question, 64 outputs are sampled. The maximum sequence length is limited to 1024, and the training batch size is also 1024. The policy model undergoes a single update after each exploration phase. DeepSeekMath-RL 7B is evaluated on benchmarks consistent with those used for DeepSeekMath-Instruct 7B. In the case of DeepSeekMath-RL 7B, GSM8K and MATH with chain-of-thought reasoning are treated as in-domain tasks, while the remaining benchmarks are considered out-of-domain tasks.

Table \ref{tab:sft_rl_math} illustrates the performance of both open- and closed-source models employing chain-of-thought and tool-enhanced reasoning on benchmarks in English and Chinese. Our observations include: 1) DeepSeekMath-RL 7B achieves accuracies of 88.2\% on GSM8K and 51.7\% on MATH using chain-of-thought reasoning. This result exceeds that of all open-source models within the 7B to 70B parameter range, as well as most closed-source models. 2) Importantly, DeepSeekMath-RL 7B is trained solely on chain-of-thought instruction tuning data from GSM8K and MATH, starting from DeepSeekMath-Instruct 7B. Despite this limited training data scope, it outperforms DeepSeekMath-Instruct 7B on every evaluation metric, demonstrating the efficacy of reinforcement learning.

\section{Discussion}
In this section, we present insights gained from our pre-training and reinforcement learning experiments.

\subsection{Lessons Learned in Pre-Training}

We begin by sharing our experiences in the pre-training phase. Unless indicated otherwise, we follow the training configurations described in Section \ref{sec:quality-policy}. It is important to note that, when mentioning the DeepSeekMath Corpus here, we refer to an 89B-token dataset collected during the second iteration of the data gathering process.

\subsubsection{Code Training Enhances Mathematical Reasoning}

A common yet unconfirmed assumption is that training on code enhances reasoning capabilities. We aim to provide a partial answer to this question, particularly within the domain of mathematics: training on code improves the model’s capacity for mathematical reasoning, both when using tools and when not.

To investigate the influence of code training on mathematical reasoning, we conducted experiments using two training paradigms:

\textbf{Two-Stage Training}

\begin{itemize}[topsep=0pt]
    \item \textbf{Code Training for 400B Tokens $\rightarrow$ Math Training for 150B Tokens}: We first train DeepSeek-LLM 1.3B on 400B code tokens, followed by training on 150B math tokens;
    \item \textbf{General Training for 400B Tokens $\rightarrow$ Math Training for 150B Tokens}: As a control, we replace code tokens with general tokens (sampled from a large-scale general corpus built by DeepSeek-AI) during the initial training phase to explore whether code tokens offer advantages over general tokens in enhancing mathematical reasoning.
\end{itemize}

\textbf{One-Stage Training}

\begin{itemize}[topsep=0pt]
    \item \textbf{Math Training for 150B Tokens}: We train DeepSeek-LLM 1.3B solely on 150B math tokens;
    \item \textbf{Training on a Mixture of 400B Code Tokens and 150B Math Tokens}: Since math training after code training leads to a decline in coding performance, we examine if combining code and math tokens in a single-stage training can still boost mathematical reasoning while alleviating catastrophic forgetting.
\end{itemize}

\paragraph{Results}
Tables \ref{tab:code-math} and \ref{tab:code-math-general-eval} present the downstream task performance under these varying training configurations.

Training on code tokens benefits program-assisted mathematical reasoning in both two-stage and one-stage training scenarios. As shown in Table \ref{tab:code-math}, within the two-stage setup, code training by itself markedly improves the capability to solve GSM8K and MATH problems using Python. Subsequent math training further enhances this ability. Notably, in the one-stage setup, mixing code and math tokens helps reduce the catastrophic forgetting observed in two-stage training, while concurrently boosting coding performance (Table \ref{tab:code-math-general-eval}) and program-assisted mathematical reasoning (Table \ref{tab:code-math}).

Code training also enhances mathematical reasoning without the use of tools. In the two-stage scenario, the initial code training phase produces moderate gains, which accelerate the following math training and ultimately achieve the highest performance. Conversely, combining code and math tokens during one-stage training impairs mathematical reasoning without tool assistance. One possible explanation is that DeepSeek-LLM 1.3B’s relatively small size limits its capacity to effectively learn both code and mathematical data simultaneously.

\subsubsection{ArXiv Papers Appear Ineffective for Enhancing Mathematical Reasoning}

ArXiv papers are frequently incorporated as part of the math pre-training dataset \citep{minerva,gpt-f,llemma,mathpile}. Nevertheless, thorough evaluations of their influence on mathematical reasoning have not been extensively explored. Contrary to what might be expected, our experiments indicate that arXiv papers do not significantly enhance mathematical reasoning. We conduct experiments on models of varying sizes, including DeepSeek-LLM 1.3B and DeepSeek-Coder-Base-v1.5 7B \citep{deepseek-coder}, using arXiv corpora processed through different pipelines:

\begin{itemize}[topsep=0pt]
    \item \textbf{MathPile} \citep{mathpile}: an 8.9B-token dataset created using cleaning and filtering heuristic rules, with over 85\% composed of scientific arXiv papers;
    \item \textbf{ArXiv-RedPajama} \citep{redpajama}: all arXiv LaTeX files with preambles, comments, macros, and bibliographies removed, amounting to 28.0B tokens.
\end{itemize}

In our trials, we individually train DeepSeek-LLM 1.3B for 150B tokens and DeepSeek-Coder-Base-v1.5 7B for 40B tokens on each arXiv dataset. The results suggest that arXiv papers do not effectively improve mathematical reasoning. When trained exclusively on arXiv datasets, both models fail to show meaningful progress or even experience declines on various mathematical benchmarks spanning different levels of difficulty used in this study. These benchmarks encompass quantitative reasoning datasets such as GSM8K and MATH (Table \ref{tab:arxiv-cot}), multiple-choice tests like MMLU-STEM (Table \ref{tab:arxiv-cot}), and formal mathematics challenges such as miniF2F (Table \ref{tab:arxiv_atp}).

Nonetheless, this finding comes with caveats and should be interpreted cautiously. We have yet to examine:

\begin{itemize}[topsep=0pt]
    \item The influence of arXiv tokens on specific math-related tasks outside the scope of this work, such as theorem informalization, which involves transforming formal statements or proofs into their informal counterparts;
    \item The effects of arXiv tokens when integrated with other data types;
    \item Whether advantages from arXiv papers might emerge at larger model scales.
\end{itemize}

Therefore, additional research is necessary, which we defer to future investigations.

\subsection{Insights of Reinforcement Learning}
\subsubsection{Towards a Unified Paradigm} In this section, we propose a unified framework to analyze various training methods, including SFT, RFT, DPO, PPO, GRPO, and conduct experiments to investigate components of this framework. Generally, the gradient with respect to the parameter $\theta$ of a training method can be expressed as:

\begin{equation}
    \nabla_{\theta}\mathcal{J}_{\textcolor{red}{\mathcal{A}}}(\theta) = \mathbb{E}[\underbrace{(q,o) \sim \textcolor{red}{\mathcal{D}}}_{Data \ Source}]\left( \frac{1}{|o|} \sum_{t=1}^{|o|}  \underbrace{GC_{{\mathcal{A}}}(q, o, t, \textcolor{red}{\pi_{{rf}}})}_{Gradient \ Coefficient}  \nabla_{\theta}\log \pi_{\theta}(o_t | q, o_{<t})\right).
\label{eq:objective}
\end{equation}

There are three principal elements involved: 1) \textit{Data Source $\mathcal{D}$}, which specifies the dataset used for training; 2) \textit{Reward Function $\pi_{{rf}}$}, which provides the reward signals utilized during training; 3) \textit{Algorithm $\mathcal{A}$}, which takes the training data and reward signals to compute the gradient coefficient $GC$, indicating the extent of penalty or reward applied to the data. We evaluate several notable techniques within this generalized framework:

\begin{itemize}
    \item  \textbf{Supervised Fine-tuning (SFT)}: SFT adjusts a pretrained model by training it on human-curated SFT datasets.
    \item \textbf{Rejection Sampling Fine-tuning (RFT)}: RFT further fine-tunes the SFT model using filtered outputs generated by sampling the SFT model on SFT questions, where outputs are selected based on answer correctness.
    \item \textbf{Direct Preference Optimization (DPO)}: DPO refines the SFT model further by fine-tuning it with augmented outputs sampled from the SFT model, employing a pairwise DPO loss function.
    \item \textbf{Online Rejection Sampling Fine-tuning (Online RFT)}: Unlike RFT, Online RFT begins with the SFT model as the initial policy and fine-tunes it on augmented outputs sampled directly from the current policy model during training.
    \item \textbf{PPO/GRPO}: PPO/GRPO also initializes the policy model with the SFT model and then reinforces it by training on outputs sampled from the real-time policy model.
\end{itemize}

We provide a summary of the elements of these approaches in Table \ref{tab:data_policy}. For a more comprehensive derivation process, please consult Appendix \ref{app:analysis-rl}.

\paragraph{Observation about Data Source} We classify the data source into two types: online sampling and offline sampling. Online sampling indicates that the training data originates from the exploration outcomes of the real-time training policy model, whereas offline sampling means the training data is derived from the sampling results of the initial SFT model. RFT and DPO adopt the offline approach, while Online RFT and GRPO utilize the online approach.

As illustrated in Figure \ref{fig:rl-analysis}, we observe that Online RFT considerably surpasses RFT on two benchmarks. Specifically, Online RFT performs similarly to RFT during the early phases of training but gains a clear advantage in later stages, showcasing the benefits of online training. This result is intuitive, as in the initial phase, the actor and the SFT model are closely aligned, causing sampled data to show only slight variations. However, in later stages, data sampled from the actor exhibits more pronounced differences, making real-time data sampling more advantageous.

\paragraph{Observation about Gradient Coefficient} The algorithm translates the reward signal into a gradient coefficient to update the model parameters. In our experiments, we categorize the reward function as either `Rule' or `Model.' The Rule approach evaluates the quality of a response based on the correctness of the answer, while the Model approach involves training a reward model to score each response. The training data for the reward model is derived from rule-based judgments. Equations \ref{eq:GC-RFT} and \ref{eq:GC-GRPO} highlight a key distinction between GRPO and Online RFT: GRPO uniquely modifies its gradient coefficient according to the reward value assigned by the reward model. This enables differentiated reinforcement and penalization of responses based on their varying scores. Conversely, Online RFT does not incorporate this mechanism; it does not penalize incorrect responses and applies uniform reinforcement to all correct responses with the same level of intensity.

As illustrated in Figure \ref{fig:rl-analysis}, GRPO outperforms online RFT, underscoring the effectiveness of modifying positive and negative gradient coefficients. Moreover, GRPO+PS achieves better results than GRPO+OS, demonstrating the advantages of employing fine-grained, step-aware gradient coefficients. Additionally, we investigate iterative RL by performing two rounds of iteration in our experiments. As depicted in Figure \ref{fig:iter-rl}, iterative RL markedly enhances performance, particularly after the first iteration.

\subsubsection{Why Does RL Work?} In this study, we apply reinforcement learning to a subset of instruction tuning data, achieving substantial improvements over the instruction tuning model. To further elucidate why reinforcement learning is effective, we assess the Pass@K and Maj@K accuracy metrics for both the Instruct and RL models across two benchmarks. As shown in Figure \ref{fig:maj-pass}, RL boosts Maj@K performance but does not improve Pass@K. This suggests that RL enhances overall model performance by making the output distribution more robust; in other words, \textbf{the improvement appears to stem from increasing the likelihood of correct responses within the TopK outputs rather than enhancing the model’s underlying capabilities.} Similarly, \citep{wang2023making} identified a \textbf{misalignment problem} in reasoning tasks within the SFT model, showing that reasoning abilities of SFT models can be improved through various preference alignment strategies \citep{yuan2023rrhf,song2023preference,wang2023making}.

\subsubsection{How to Achieve More Effective RL?}
\label{sec:effec_rl}
We show that RL performs particularly well on mathematical reasoning tasks and present a unified framework to understand several representative training methods. Within this framework, all methods can be viewed as either direct or simplified RL approaches. As summarized in Equation \ref{eq:objective}, there are three main components: Data Source, Algorithm, and Reward Function. We outline some prospective directions related to these components.

\paragraph{Data Source} The data source serves as the fundamental input for all training methods. In the RL context, the data source specifically refers to unlabeled questions with outputs generated by sampling from the policy model. In this work, we only use questions from the instruction tuning stage and employ a basic nucleus sampling technique for output generation. We believe this may explain why our RL pipeline improves only Maj@K performance. Going forward, we plan to explore our RL pipeline using out-of-distribution question prompts combined with \textbf{advanced sampling (decoding) strategies}, such as those based on tree-search methods \citep{yao2023tree}. Furthermore, \textbf{efficient inference techniques} \citep{xia-etal-2023-speculative,leviathan2023fast,kwon2023efficient,xia2024unlocking}, which influence the exploration efficiency of policy models, will also play a critically important role.

\paragraph{Algorithms} Algorithms utilize the data and reward signals to compute the gradient coefficient, which is then used to update the model parameters. According to Equation \ref{eq:objective}, current approaches largely \textbf{TRUST} the reward function's signal to either increase or decrease the conditional probability of a given token. Nevertheless, it is unrealistic to guarantee that the reward signal is consistently reliable, especially in highly complex tasks. For instance, even the PRM800K datasets \citep{lightman2023let}, carefully annotated by skilled annotators, contain roughly 20\% erroneous annotations\footnote{\url{https://github.com/openai/prm800k/issues/12\#issuecomment-1728491852}}. Therefore, we intend to investigate reinforcement learning algorithms that exhibit robustness to noisy reward signals. We anticipate that such \textbf{WEAK-TO-STRONG} \citep{burns2023weak} alignment methods will fundamentally transform learning algorithms.

\paragraph{Reward Function} The reward function serves as the origin of the training signal. In reinforcement learning, the reward function is commonly represented by a neural reward model. We identify three critical directions for reward models: 1) \textbf{Improving the generalization capability of the reward model.} It is essential for the reward model to generalize effectively to handle out-of-distribution queries and complex decoding outputs; otherwise, reinforcement learning might only maintain the existing distribution of LLMs without enhancing their core abilities; 2) \textbf{Capturing the uncertainty of the reward model.} This uncertainty could serve as a crucial link between weak reward models and weak-to-strong learning algorithms; 3) \textbf{Developing efficient, high-quality process reward models} that provide detailed training signals for the reasoning process \citep{lightman2023let,wang2023math}.

\section{Conclusion, Limitation, and Future Work}  
We introduce DeepSeekMath, which surpasses all open-source models on the competitive MATH benchmark and nears the performance level of proprietary models. DeepSeekMath~starts from DeepSeek-Coder-v1.5 7B and is continually trained over 500B tokens, with a substantial portion—120B tokens—derived from mathematical content collected via Common Crawl. Our comprehensive ablation study reveals that web pages hold great promise for providing high-quality mathematical data, whereas arXiv contributes less advantage than initially anticipated. We propose Group Relative Policy Optimization (GRPO), a variant of Proximal Policy Optimization (PPO), which significantly enhances mathematical reasoning abilities while consuming less memory. Experimental results demonstrate that GRPO remains effective even when DeepSeekMath-Instruct 7B attains high benchmark scores. Moreover, we offer a unified framework to interpret various approaches and highlight multiple avenues for advancing reinforcement learning methods.

Despite DeepSeekMath's strong performance on quantitative reasoning tasks, its proficiency in geometry and theorem-proving lags behind closed models. For example, during our preliminary testing, the model struggled with problems involving triangles and ellipses, suggesting possible biases in data selection during pre-training and fine-tuning phases. Furthermore, limited by model size, DeepSeekMath~underperforms compared to GPT-4 in few-shot learning scenarios. While GPT-4 benefits from few-shot prompts, DeepSeekMath~exhibits comparable outcomes between zero-shot and few-shot evaluations. Moving forward, we aim to enhance our curated data selection process to build a higher-quality pre-training corpus. Additionally, we intend to investigate the promising strategies outlined in Section \ref{sec:effec_rl} to further advance reinforcement learning techniques for large language models.

\end{document}